\subsection{破碎}

定义 10. --- 称 T 对 \(\mu\) 的破碎，或 \(\mu\) -破碎，为 T 的两两不相交紧子集的任一局部可数族 \((K_{\alpha})_{\alpha \in A}\) ，如果集 \(N = T - \bigcup_{\alpha} K_{\alpha}\) 局部 \(\mu\) -可忽略。

命题 9. --- a) 存在一个破碎 \((\mathrm{K}_{\alpha})_{\alpha \in \mathrm{A}}\) ，它是 T 对 \(\mu\) 的破碎。 b) 设 \((\mathrm{K}_{\alpha})_{\alpha \in \mathrm{A}}\) 是 T 对 \(\mu\) 的一个破碎。若 \(\mu_{\alpha}\) 是由 \(\mu_{\mathrm{K}_{\alpha}}\) 定义的 T 上的测度 (No.~3, 例 2)，则族 \((\mu_{\alpha})_{\alpha \in \mathrm{A}}\) 可和，其和等于 \(\mu\) ，且对每个函数 \(f \in \mathcal{F}_{+}(\mathrm{T})\) ，

\[
\mu^ {\bullet} (f) = \sum_ {\alpha \in A} \mu_ {\alpha} ^ {\bullet} (f) = \sum_ {\alpha \in A} \mu_ {K _ {\alpha}} ^ {\bullet} (f _ {K _ {\alpha}}). ^ {(1)}\tag{7}
\]

为使 T 到拓扑空间 G（Hausdorff 与否均可）中的映射 g 是 \(\mu\) -可测的，必须且只需 \(g_{K_{\alpha}}\) 是 \(\mu_{K_{\alpha}}\) -可测的（对每个 \(\alpha \in A\)）。

A) 破碎的存在性：

本证明是 Ch. IV, §5, No.~9 的 Prop. 14 的证明稍作修改后的重复。设 \(\mathcal{K}\) 是紧子集 \(K\) ⊂ \(T\) 且 \(\operatorname{Supp}(\mu_K) = K\) 者的全体，\(\mathcal{H}\) 是子集 \(\mathcal{L}\) ⊂ \(\mathcal{K}\) （由两两不相交的集组成，按包含关系排序）的全体。首先证明每个元素 \(\mathcal{L}\) ∈ \(\mathcal{H}\) 都局部可数。设 \(x\) 是 \(T\) 的一点，\(V\) 是 \(x\) 的开邻域且 \(\mu^\bullet(V) < +\infty\) ；设 \(\mathfrak{L}_V\) 是 \(K \in \mathfrak{L}\) 中与 \(V\) 相交者的全体。若 \((K_i)_{1 \leqslant i \leqslant n}\) 是 \(\mathfrak{L}_V\) 中互异元素的一个有限序列，则由 Prop. 4 的推论，

\[
\sum_ {i = 1} ^ {n} \mu^ {\bullet} (\mathrm{K} _ {i} \cap \mathrm{V}) = \mu^ {\bullet} \left(\mathrm{V} \cap \left(\bigcup_ {i = 1} ^ {n} \mathrm{K} _ {i}\right)\right) \leqslant \mu^ {\bullet} (\mathrm{V}),
\]

因为诸 \(K_{i}\) 两两不相交。于是，

\[
\sum_ {\mathrm{K} \in \mathfrak {L} _ {\mathrm{V}}} \mu^ {\bullet} (\mathrm{K} \cap \mathrm{V}) <   + \infty .
\]

而 \(\mu^{\bullet}(\mathbf{K}\cap \mathbf{V}) = \mu_{\mathbf{K}}^{\bullet}(\mathbf{K}\cap \mathbf{V}) > 0\) 对每个 \(\mathbf{K}\in \mathfrak{L}_{\mathbf{V}}\) 成立，因为 \(\mathbf{K}\cap \mathbf{V}\) 非空且在 \(\mathbf{K}\) 中开，而 \(\mu_{\mathbf{K}}\) 的支集是整个 \(\mathbf{K}\) ；因此 \(\mathfrak{L}_{\mathbf{V}}\) 可数，从而 \(\mathfrak{L}\) 确实局部可数。显然 \(\mathcal{H}\) 是归纳的且非空（有 \(\varnothing \in \mathcal{H}\) ）。于是设 \(\mathfrak{H}\) 是 \(\mathcal{H}\) 的一个极大元素。我们将证明集 \(\mathrm{N} = \mathrm{T} - \bigcup_{\mathrm{K}\in \mathfrak{H}}\mathrm{K}\)

局部可忽略。由 Prop. 2，只需验证对每个紧集 L 有 \(\mu^{\bullet}(\mathbf{N}\cap\mathbf{L})=0\) ，或者等价地 \(\mu_{\mathrm{L}}^{\bullet}(\mathrm{N}\cap\mathrm{L})=0\) 。我们用反证法。于是设 \(\mu_{\mathrm{L}}^{\bullet}(\mathrm{N}\cap\mathrm{L})>0\) 。由于与 L 相交的 \(\mathbf{K}\in\mathfrak{H}\) 全体可数，N∩L 是 \(\mu_{\mathrm{L}}\) -可测的；因此存在含于 N∩L 的紧集 J 使 \(\mu_{\mathrm{L}}^{\bullet}(\mathrm{J})>0\) 。设 S 是非零测度 \((\mu_{\mathrm{L}})_{\mathrm{J}}=\mu_{\mathrm{J}}\) 的支集；它含于 N，测度 \(\mu_{\mathrm{S}}\) 非零，且 \(\operatorname{Supp}(\mu_{\mathrm{S}})=\mathrm{S}\) (Ch. IV, §5, No.~7, Lemma 2)。于是集 \(\mathfrak{H}\cup\{\mathrm{S}\}\) 属于 \(\mathscr{H}\) ，与 \(\mathfrak{H}\) 的极大性矛盾。这就证明了破碎的存在性。

\begin{enumerate}
\def\labelenumi{\Alph{enumi})}
\setcounter{enumi}{1}
\tightlist
\item
  (7) 的证明：
\end{enumerate}

对每个 \(\alpha \in \mathbf{A}\) ，由 No.~3 的公式 (3) 及 No.~2 的 Prop. 2 a)，有 \(\mu_{\alpha}^{\bullet}(f) = \mu_{\mathrm{K}_{\alpha}}^{\bullet}(f_{\mathrm{K}_{\alpha}}) = \mu^{\bullet}(f\varphi_{\mathrm{K}_{\alpha}})\) ；这些公式表明负担 \(\sum_{\alpha \in \mathbf{A}} \mu_{\alpha}^{\bullet}\) \(\leqslant \mu^{\bullet}\) ，从而族 \((\mu_{\alpha})_{\alpha \in \mathbf{A}}\) 可和 (No.~7, Prop. 7)。于是只需证明 \(\mu = \sum_{\alpha \in \mathbf{A}} \mu_{\alpha}\) ，即建立公式

\[
\mu_ {\mathrm{K}} ^ {\bullet} = \sum_ {\alpha \in \mathrm{A}} (\mu_ {\alpha}) _ {\mathrm{K}} ^ {\bullet}\tag{8}
\]

对 T 的每个紧子集 K 成立。现固定 K，由 \(\alpha \in A\) 中使 \(K_{\alpha}\) 与 K 相交者所成的集 A' 可数。设 \(g \in \mathcal{F}_{+}(K)\) ；则 \(g^{0} = g^{0}\varphi_{N} + \sum_{\alpha \in A} g^{0}\varphi_{K_{\alpha}}\) ，且 \(g^{0}\varphi_{K_{\alpha}} = 0\) （对 \(\alpha \in A - A'\)）；由 No.~5 的 Prop. 4 推出 \(\mu^{\bullet}(g^{0}) = \sum_{\alpha \in A} \mu^{\bullet}(g^{0}\varphi_{K_{\alpha}})\) ，由此

\[
\mu_ {\mathrm{K}} ^ {\bullet} (g) = \mu^ {\bullet} (g ^ {0}) = \sum_ {\alpha \in \mathrm{A}} \mu^ {\bullet} (g ^ {0} \varphi_ {\mathrm{K} _ {\alpha}}) = \sum_ {\alpha \in \mathrm{A}} \mu_ {\alpha} ^ {\bullet} (g ^ {0}) = \sum_ {\alpha \in \mathrm{A}} (\mu_ {\alpha}) _ {\mathrm{K}} ^ {\bullet} (g);
\]

于是 (8) 得证。

C) 可测性：

为使定义在 T 上的函数 g 是 \(\mu\) -可测的，必须且只需它是 \(\mu_{\alpha}\) -可测的（对每个 \(\alpha \in A\)）(No.~7, Prop. 8)；而这等于说 \(g_{K_{\alpha}}\) 是 \(\mu_{K_{\alpha}}\) -可测的（对所有 \(\alpha \in A\)）(No.~5, 例)。证毕。

如同 Ch. IV, §5, No.~9 的 Prop. 14，可以要求诸紧集 \(K_{\alpha}\) 属于预先给定的 T 的紧子集的一个 \(\mu\) -稠密集。我们只需要下面这个将直接证明的结果：

命题 10. --- 若 g 是以拓扑空间 G（Hausdorff 与否均可）为值的 \(\mu\) -可测映射，则存在 T 的一个 \(\mu\) -破碎 \((\mathbf{L}_{\beta})_{\beta\in\mathbf{B}}\) ，使限制 \(g_{L_{\beta}}\) （对所有 \(\beta\in B\)）连续。

取一个破碎 \((\mathrm{K}_{\alpha})_{\alpha\in\mathrm{A}}\) （T 对 \(\mu\) 的破碎）。由于映射 g 可测，对每个 \(\alpha\in A\)，存在把 \(K_{\alpha}\) 分解为一列紧集 \((\mathrm{K}_{\alpha n})\) 与一个局部可忽略集 \(N_{\alpha}\) 的划分，使 g 到每个集 \(K_{\alpha n}\) 的限制连续。于是族 \((\mathrm{K}_{\alpha n})_{(\alpha,n)\in\mathrm{A}\times\mathbf{N}}\) 就是所求的破碎。因为它是局部可数的，且集 \(N' = N \cup \left( \bigcup_{\alpha} N_{\alpha} \right)\) 局部可忽略，这是由于一个紧集至多与可数无穷多个集 \(N_{\alpha}\) 相交。

附论. --- 设 \((\mathrm{K}_{\alpha})_{\alpha\in\mathrm{A}}\) 是 T 的一个破碎，令 \(N=T-\bigcup_{\alpha}K_{\alpha}\) 。用 \(T'\) 表示在 T 上赋予诸子空间 \(K_{\alpha}\) 的拓扑之和拓扑以及 N 上任一局部紧拓扑所得到的局部紧空间（除非特别说明相反情形，N 总是赋予离散拓扑）。对每个 \(\alpha \in A\) ，设 \(i_{\alpha}\) 是 \(K_{\alpha}\) 到 \(T'\) 中的典范单射，\(\mu_{\alpha}'\) 是 \(T'\) 上 \(\mu_{K_{\alpha}}\) 在 \(i_{\alpha}\) 下的像测度。族 \((\mu_{\alpha}')\) 可和：因为若 f 是 \(T'\) 上具紧支集的连续函数，则 Supp(f) 与 \(K_{\alpha}\) 只对有限个指标 \(\alpha\) 相交。令 \(\mu' = \sum_{\alpha \in A} \mu_{\alpha}'\) 。

由于集 N 对 \(\mu'\) 局部可忽略——因为它对每个 \(\mu'_{\alpha}\) 如此 (Prop. 9)——故族 \((\mathrm{K}_{\alpha})_{\alpha \in \mathrm{A}}\) 是一个 \(\mu'\) -破碎（对 \(\mathrm{T}'\) 而言）；而 \(\mu'\) 在 \(\mathrm{K}_{\alpha}\) 上诱导的测度显然是 \(\mu_{\mathrm{K}_{\alpha}}\)，公式 (7) 应用于 \(\mu\) 与 \(\mu'\) 表明 \(\mu^{\bullet} = \mu'^{\bullet}\) 。类似地，Prop. 9 陈述的最后一条应用于 \(\mu\) 与 \(\mu'\) 表明，两个测度 \(\mu\) 与 \(\mu'\) 的可测映射是相同的。

这两条性质使我们能把关于 \(\mu\) 的积分理论的几乎全部内容化归为对局部紧空间已建立的理论。这些考虑将在 No.~10 中展开。

下面是破碎概念的另一个应用：

命题 11. --- 设 X 是 T 的一个 \(\mu\) -可测子集。存在 X 的紧子集的一个局部可数族 \((\mathrm{L}_{\alpha})_{\alpha \in \mathbf{A}}\) ，两两不相交，使 \(\mathrm{X} - \bigcup_{\alpha \in \mathbf{A}} \mathrm{L}_{\alpha}\) 局部 \(\mu\) -可忽略。此外，若 X 是一列可测集 \((\mathrm{X}_n)\) （满足 \(\mu^{\bullet}(\mathrm{X}_n) < +\infty\)）之并，则 \(\alpha \in \mathbf{A}\) 中使 \(\mu^{\bullet}(\mathrm{L}_{\alpha}) \neq 0\) 者所成的集 B 可数，且 \(\mathrm{X} - \bigcup_{\beta \in \mathbf{B}} \mathrm{L}_{\beta}\) 局部 \(\mu\) -可忽略。

设 \(f\) 是 \(X\) 的特征函数，\((K_{\alpha})_{\alpha \in A}\) 是 \(T\) 的一个破碎，使 \(f\) 到每个 \(K_{\alpha}\) 的限制连续 (Prop. 10)。于是集 \(L_{\alpha} = K_{\alpha} \cap X\) 对每个 \(\alpha \in A\) 是紧的，而 \((L_{\alpha})_{\alpha \in A}\) 即为所求的族。转到第二条断言；显然可设诸可测集 \(X_{n}\) 不相交，且只需对其中每一个证明该断言。换句话说，必要时改变记号，可设 \(\mu^{\bullet}(X) < +\infty\) 。于是集 \(B\) ——即 \(\alpha \in A\) 中使 \(\mu^{\bullet}(L_{\alpha}) > 0\) 者——可数，剩下只要证明集 \(N = \bigcup_{\alpha \in A - B} L_{\alpha}\) 局部可忽略。设 \(K\) 是一个紧集；族

\((\mathrm{L}_{\alpha})_{\alpha\in\mathrm{A}}\) 局部可数，故集 \(\mathrm{K}\cap\mathrm{N}\) 是族 \((\mathrm{K}\cap\mathrm{L}_{\alpha})_{\alpha\in\mathrm{A}-\mathrm{B}}\) 的一个可数子族之并，从而该集局部可忽略。于是 \(\mathrm{N}\) 亦然 (No.~2, Prop. 2)，命题得证。

